\subsection{A certified range of sensing angles}\label{sec:sensing_angle}
\begin{corollary}[Non-aligned risk separation]\label{cor:sensing_angle}
Use the Gaussian model of Theorem~\ref{thm:ind_noise} and its
template and mass box with $\alpha=10^{-3}$. In the same template
coordinates, take $B_\varphi=R_\varphi\operatorname{diag}(\alpha,1)
R_\varphi^\top=A^\top A$, with $R_\varphi$ a planar rotation.
For exact RED from $x_0=A^\top y$, take $\lambda=1$ and $v_L=S$.
As $\tau,\nu\to0$ independently, the restricted risk has limiting
lower bound $.147$ uniformly over all finite budgets and histories,
and the $0/25/0$ preset achieves $O(S^2)$ in $25+n$ calls with
$n$ as in~\eqref{eq:fixed_geom}. Both bounds hold uniformly for
$|\varphi|\le .020$ radians (about $1.15^\circ$), with the scale sets
of Theorem~\ref{thm:ind_noise}.
\end{corollary}

For $e=b-B_\varphi\mu_Z$ compare with the aligned update at
$b_0=B_0\mu_Z+e$. The update difference is
$-\gamma(B_\varphi-B_0)(x-\mu_Z)$.
The diagonal entries of $B_\varphi-B_0$ have magnitude at most
$\varphi^2$ and the off-diagonal entries at most $|\varphi|$.
Thus the perturbation can be bounded separately on each face.
All rotated maps are nondecreasing in each output's own input
coordinate, since $B_{\varphi,jj}\le1$ and
$\partial_j\bar\mu_{v,j}\ge0$. These face tests imply invariance
of the respective rectangles.

\paragraph{Common barrier.}
For class 1, the backprojection starts in
$W=[-.5,.3]\times[-.027,.51]$ on $E$.
On the right face, let $M$ bound $\bar\mu_{v,1}$ from above
and $U$ bound $|t|$ on the part of the face under consideration.
With $f=.020$, the displacement is at most
\[
 \gamma[.00095+.999(M-.3)+.7f^2+fU].
\]
At low scales use $(M,U)=(.287729,.51)$. At high scales,
the negative competitor's odds increase with $t$.
Splitting at $t=.15$ yields $(M,U)=(.29300,.15)$ and
$(.25984,.51)$ on the lower and upper segments.
All three bounds give an inward margin of more than $.00020$.
The other three faces
have margins above $.047$, $.0047$, and $.0030$ at the left,
bottom, and top, respectively (Supplement S27). Hence $s\le.3$
for any history using the restricted union.

\paragraph{Entry and terminal risk.}
The original observation classifier works on $E$. Class 1 uses
a widened transverse strip $[-.027,.06]$. The uniform slope
bound is $.564$ and the entrance surplus exceeds $.01563$.
With horizontal rotation correction below $.000402$ per
update, the lower comparison is
\begin{align*}
 \underline s_{k+1}&=\underline s_k+
    \gamma(.999)[.01563+.436(.53-\underline s_k)]-.000402,\\
 \underline s_0&=.00075,\qquad \underline s_{25}>.5306.
\end{align*}
The prepared region remains invariant through the switch.
Classes 0 and 2 start in their terminal rectangles, with the
class-2 right face widened to $s=.05$. The terminal
distance gaps are still at least $.06$.
For the actual center
$m_\varphi=(B_\varphi+c_LI)^{-1}(b+c_L\mu_Z)$, the
residual recursion has linear factor at most $.751$ and
posterior leakage $O(e^{-.03/v_L})$.
The noise covariance satisfies
$\operatorname{Cov}(e)=\tau^2B_\varphi^2+\nu^2B_\varphi\preceq S^2I$.
The resulting Gaussian tail and terminal estimate give the
$O(S^2)$ upper bound uniformly in the angle range. The common
barrier retains the $.147$ limiting lower bound.
